% SPDX-License-Identifier: Apache-2.0
\section{An Immediate Decoder}
\label{sec:x-imm}

The bit operations a decoder is made of: \code{slice::<LO, LEN>}
takes a field out of a word, \code{concat::<K, M>} puts one value
above another with the result width stated, and \code{sext::<M>} and
\code{zext::<M>} extend by the top bit or by zeros. Each lowers to the
form the target language has, a part select, a concatenation, a
replication of the top bit in Verilog and \code{resize} in VHDL. The
five immediates of RV32I are those operations as the manual draws
them, and \code{select!} chooses the one the opcode calls for.

\code{select!} is \code{match} on a value: the first pattern that
matches chooses, a pattern may carry an \code{if} guard, and the last
arm is the default. The name says what the hardware is, a chain of
multiplexers in priority order, which is what the lowering makes of
it; \code{case!} is the same choice among drives, \code{select!} the
same choice among values. Both take an integer literal as a pattern,
so an opcode can be matched as the manual writes it. The decoder is
checked under nvc against its trace, Section~\ref{sec:x-checked}, and
it is the decode of the Vreteno core, lifted out.

\begin{figure}[htbp]
\centering
\input{imm_timing.tex}
\caption{The decoder: a word of each format, and its immediate and
destination a cycle later.}
\label{fig:x-imm}
\end{figure}

\lstinputlisting[caption={\code{ex\_imm.rs}. Run with \code{bazel run //lib/examples:ex\_imm}.},label={lst:x-imm}]{lib/examples/ex_imm.rs}

\lstinputlisting[style=ascii,caption={What \code{ex\_imm} printed when the build ran it: the words, then the lowering.},label={lst:x-imm-out}]{out_imm.txt}
